\chapter{Visão Geral}

O microserviço \textit{PlantDiseaseAI} atua como um serviço especializado dentro de um ecossistema maior. Sua função não é concentrar todas as responsabilidades do sistema agrícola, mas servir como unidade independente responsável pela inferência futura de doenças em plantas. A visão geral da solução pode ser compreendida como uma sequência de etapas: uma API recebe a solicitação, grava o estado inicial, armazena a imagem, publica uma mensagem assíncrona e, posteriormente, um consumidor aciona o processamento do pipeline de IA.

A comunicação com outros serviços ocorre por dois mecanismos complementares. O primeiro é HTTP, utilizado pela API para expor a criação e a consulta do estado de uma análise. O segundo é mensageria assíncrona, utilizada para desacoplar o início da análise do processamento computacionalmente pesado. Ao concluir o pipeline, o serviço retorna o resultado ao sistema solicitante por callback HTTP, preservando a interoperabilidade com backends externos.

Essa abordagem permite que o sistema externo não precise conhecer a complexidade interna do processamento. Ele apenas envia uma imagem e um contexto mínimo da requisição, acompanha o progresso e recebe o resultado final. Em ambientes de produção, isso reduz a carga sobre a camada síncrona e cria uma base mais robusta para escalabilidade horizontal e tolerância a falhas.

\section{Responsabilidades principais}

As responsabilidades do microserviço incluem o registro da solicitação de análise, a persistência temporária do job, o armazenamento da imagem, o disparo de eventos, a execução do pipeline de IA e a preparação do resultado final para callback. Cada responsabilidade está confinada a uma camada específica para manter clareza arquitetural e facilitar manutenção.

\section{Fluxo de alto nível}

O fluxo lógico da solução é composto pelas seguintes etapas:
\begin{enumerate}
    \item o cliente chama a API para criar uma nova análise;
    \item o serviço registra o job e armazena a imagem;
    \item uma mensagem de solicitação é publicada;
    \item o consumidor recebe a mensagem;
    \item o processor carrega o job e atualiza o status;
    \item o pipeline executa pré-processamento, carregamento do modelo, execução e pós-processamento;
    \item o resultado é persistido e enviado ao sistema solicitante.
\end{enumerate}

\begin{figure}[H]
    \centering
    \fbox{\parbox{0.85\textwidth}{\centering Espaço reservado para o diagrama geral da arquitetura do microserviço PlantDiseaseAI.}}
    \caption{Visão geral da arquitetura do microserviço}
    \label{fig:visao-geral-arquitetura}
\end{figure}

\section{Integrações externas}

A solução se comunica com um backend externo responsável por iniciar a análise e consumir os callbacks, com RabbitMQ para mensageria e, futuramente, com o modelo treinado de IA. O desenho atual já antecipa a substituição do pipeline placeholder por uma implementação real, sem alterar a API pública do serviço.
